% Part: sets-functions-relations
% Chapter: functions
% Section: functions-relations

\documentclass[../../../include/open-logic-section]{subfiles}

\begin{document}

\olfileid{sfr}{fun}{rel}

\olsection{సంబంధాలుగా ప్రమేయాలు}

\begin{explain}
$A$లోని \tecase{acc}{!!{element}s}, $B$లోని \tecase{dat}{!!{element}s}
ప్రతిచిత్రించే ప్రమేయం $A$, $B$ల మధ్య ఒక సంబంధాన్ని నిర్వచిస్తుందనేది
స్పష్టమే: $f(x) = y$ అయినప్పుడూ, అప్పుడే, $x$, $y$ల మధ్య ఉండే
సంబంధమే అది. నిజానికి, ఉదాహరణకు గణితపు నిర్మాణ భాగాలను తగ్గించాలనుకుంటే,
$f$ ప్రమేయాన్ని ఈ సంబంధంతోనే, అంటే ఒక జతల సమితితోనే,
\emph{ఒకటిగా పరిగణించవచ్చు}. అప్పుడు ఒక ప్రశ్న తలెత్తుతుంది:
ఈ విధంగా ప్రమేయాలను నిర్వచించే సంబంధాలు ఏవి?
\end{explain}

\begin{defn}[ప్రమేయపు గ్రాఫు]
$f\colon A \to B$ ఒక ప్రమేయం అనుకుందాం.
$f$ యొక్క \emph{గ్రాఫు} అంటే కింది విధంగా నిర్వచించే
$R_f \subseteq A \times B$ అనే సంబంధం:
\[
R_f = \Setabs{\tuple{x,y}}{f(x) = y}.
\]
\end{defn}

\begin{explain}
మూలకాధారిత సమానత్వం ప్రకారం ప్రమేయపు గ్రాఫు ఏకైకంగా నిర్ణయించబడుతుంది.
అంతేకాక, సమితులకు వర్తించే మూలకాధారిత సమానత్వం, ప్రమేయాల విషయంలో
అంతర్లీనంగా వాడే సమానత్వ సూత్రాన్ని వెంటనే సమర్థిస్తుంది:
$f$, $g$లకు ఒకే ప్రవేశమూ సహప్రవేశమూ ఉంటే, అన్ని విలువల వద్ద అవి
ఏకీభవించినప్పుడు అవి ఒకటే అవుతాయి.

అదేవిధంగా, ఒక సంబంధం ``ప్రమేయాత్మకంగా'' ఉంటే, అది ఒక ప్రమేయపు గ్రాఫు.
\end{explain}

\begin{prop}\ollabel{prop:graph-function}
$R \subseteq A \times B$ కింది షరతులను తీరుస్తుందని అనుకుందాం:
\begin{enumerate}
\item $Rxy$, $Rxz$ అయితే $y = z$; అలాగే
\item ప్రతి $x \in A$కూ $\tuple{x,y} \in R$ అయ్యే $y \in B$ ఏదో ఒకటి ఉంటుంది.
\end{enumerate}
అప్పుడు $R$ అనేది $Rxy$ అయినప్పుడూ, అప్పుడే, $f(x) = y$ అయ్యేలా
నిర్వచించిన $f\colon A \to B$ అనే ప్రమేయపు గ్రాఫు.
\end{prop}

\begin{proof}
$Rxy$ అయ్యే $y$ ఉందనుకుందాం.
$Rxz$ అయ్యే మరో $z \neq y$ ఉంటే, $R$పై విధించిన షరతు ఉల్లంఘించబడుతుంది.
కాబట్టి $Rxy$ అయ్యే $y$ ఉంటే, ఆ $y$ ఏకైకం; అందువల్ల $f$ సునిర్వచితం.
$R_f = R$ అనేది స్పష్టమే.
\end{proof}

\begin{explain}
ప్రతి ప్రమేయం $f\colon A \to B$కూ ఒక గ్రాఫు ఉంటుంది; అంటే
$f(x) = y$ ద్వారా నిర్వచించే $A \times B$పై ఒక సంబంధం ఉంటుంది.
మరోవైపు, \olref{prop:graph-function}లో పేర్కొన్న ధర్మాలు గల ప్రతి
$R \subseteq A \times B$ సంబంధమూ ఒక ప్రమేయం $f \colon A \to B$ యొక్క గ్రాఫు.
ప్రమేయాలకూ వాటి గ్రాఫులకూ గల ఈ సన్నిహిత సంబంధం వల్ల,
ప్రమేయాన్ని కేవలం దాని గ్రాఫుగానే భావించవచ్చు. మరో మాటలో,
ప్రమేయాలను కొన్ని నిర్దిష్ట సంబంధాలతో, అంటే కొన్ని క్రమిత బహుళకాల
సమితులతో, ఒకటిగా పరిగణించవచ్చు.
\oliflabeldef{sfr:rel:ref:sec}{అయితే ఈ ``ఒకటిగా పరిగణించడం'' యొక్క
ఉద్దేశం \olref[sfr][rel][ref]{sec}లో ఉన్నదేనని గమనించండి:
ఇది ప్రమేయాల అధిభౌతిక స్వరూపం గురించి వాదన కాదు; ప్రమేయాలను కొన్ని
సమితులుగా \emph{పరిగణించడం} సౌకర్యవంతమన్న పరిశీలన మాత్రమే.
ఈ సౌలభ్యానికి ఒక కారణమిది: }{}ఇప్పుడు సంబంధాలపై చేసినవంటి ప్రక్రియలనే
ప్రమేయాలపై చేయడాన్ని పరిగణించవచ్చు
(\olref[sfr][rel][ops]{sec} చూడండి). ముఖ్యంగా:
\end{explain}

\begin{defn}\ollabel{defn:funimage}
$C\subseteq A$ కాగా, $f \colon A \to B$ ఒక ప్రమేయం అనుకుందాం.

$f$ను $C$కు \emph{పరిమితం చేయగా} వచ్చే ప్రమేయం
$\funrestrictionto{f}{C}\colon C \to B$:
ప్రతి $x \in C$కూ $(\funrestrictionto{f}{C})(x) = f(x)$గా దాన్ని నిర్వచిస్తాం.
మరో మాటలో, $\funrestrictionto{f}{C} = \Setabs{\tuple{x, y} \in R_f}{x \in C}$.

$f$ను $C$కు \emph{వర్తింపజేసిన ఫలితం}
$\funimage{f}{C} = \Setabs{f(x)}{x \in C}$.
దీన్ని $f$ కింద $C$ యొక్క \emph{ప్రతిబింబం} అని కూడా అంటాం.
\end{defn}

\begin{explain}
ఈ నిర్వచనాల ప్రకారం, ఏ ప్రమేయం $f$కైనా
$\ran{f} = \funimage{f}{\dom{f}}$ అవుతుంది.
\oliflabeldef{sfr:rel:ops:sec}{\olref[sfr][rel][ops]{sec}లోని నిర్వచనాలను,
ప్రమేయాలను సంబంధాలతో ఒకటిగా పరిగణించడాన్ని దృష్టిలో పెట్టుకుంటే,
ఈ భావనలు ఆశించినట్లే ఉంటాయి. అయితే మరో రెండు ప్రక్రియలకు---విలోమాలు,
సాపేక్ష లబ్ధాలు---కొంత అదనపు వివరణ అవసరం.
దాన్ని \olref[inv]{sec}, \olref[cmp]{sec}లలో ఇస్తాం.}{}
\end{explain}

\end{document}
